% Verbatim copy of /Users/alingalatan/workspaces/grothe/handoff/manuscript/legacy/ch33-gpt-polarized.tex (file dated 2026-10-25). Only this line was added.
% Consolidated source module (formerly Chapter 33). Owner: lower-bound/08-shared-structure.tex
\section{Lower-bound GPT: exact K for the whole Gaussian covariance class of the odd and even sectors (20 October)}
\label{ch:polarized-profile-certificate}
\chapsummary{Exact K for the matched Gaussian covariance class, with all stated orientations and penalties. This is a class theorem, not a replacement of a general cubic by a matched Gaussian.}

\paragraph{Provenance and status (IDE note, 20 October).} Lower-bound GPT's \code{proof/ch32-polarized-profile.tex} (renumbered: our \autoref{ch:stronger} took the number 32) from
\pathref{gpt-reviews/polarized-profile-20261020/}, ported verbatim apart from this header and the closing reading note. \emph{Our replay} \st{C}: \code{python3 -I verify\_all.py} (standard library
and Apple clang 17, \code{\_\_int128}) from an unrelated directory: \textbf{PASS in 300\,s}. Both parameter trees were rebuilt from regenerated rational data: the $A/B$ tree (631
numerical, 17 analytic and 5 tail leaves, 1{,}226 vertices) and the $A/B$ plus improper-even tree (4{,}175 numerical, 66 analytic and 15 tail leaves, 9{,}685 vertices). The worst numerical
upper is $96683362167072/2^{46}$, below the independent benchmark lower bound $96683367193980/2^{46}\le S$. Log \pathref{gpt-reviews/polarized-profile-20261020/replay\_ide.log}.

\providecommand{\eps}{\varepsilon}
\providecommand{\E}{\mathbb E}
% Proposed addition; live manuscript was not modified.
% Numerical bounds on the Grothendieck constant are unchanged.
\subsection{Result and precise scope}
Set
\[
 a_0=\frac{699}{1250},\quad b_0=\frac{179}{1250},\quad
 \eps=\frac{1891}{2500},\quad
 \ell(t)=\sqrt t(a_0+b_0t),\quad A(t)=t-1,\quad B(t)=t/\sqrt2.
\]
Let $T$ be exponential of mean one and $G$ a standard proper complex
Gaussian, independent of $T$. Write
\[
 S=\sup_{v\ge0}\left\{2\E|\ell(T)+\sqrt vG|-v/\eps\right\}.
\]
The main result closes the \emph{Gaussian covariance layer} of the odd
sectors. It does not prove Conjecture K for general non-Gaussian shared
cubic drivers, and it does not change the certified Grothendieck bound.

\begin{theorem}[All amplitudes, all covariance orientations]\label{polar:main}
Fix $g,\alpha,\beta\ge0$. Suppose that, conditionally on $T=t$,
$R$ is a centered real two-dimensional Gaussian, possibly degenerate,
with complex variance and pseudovariance
\[
 V(t)=g^2+\alpha^2 A(t)^2+\beta^2 B(t)^2,
 \qquad |P(t)|\le 2\alpha\beta |A(t)B(t)|.
\]
The argument of $P(t)$ may be an arbitrary measurable function of $t$.
Then
\[
 \boxed{\quad
 2\E|\ell(T)+R|-\frac{g^2+\alpha^2+\beta^2}{\eps}
 \le S-\frac{\alpha^2}{10}-\frac{\beta^2}{4}.
 \quad}                                                   \tag{1}
\]
There are no restrictions on amplitudes or on the number of Gaussian
drivers. This is a proved mathematical statement with a finite exact
arithmetic certificate, not a Lean formalization.
\end{theorem}

For example, (1) applies to
\[
 R=gG+A(T)\sum_j u_jX_j+B(T)\sum_jv_j\overline{X_j},
 \qquad \alpha=\|u\|_2,\quad\beta=\|v\|_2,
\]
where $G,X_1,X_2,\ldots$ are independent proper Gaussians. Thus shared
A/B drivers, every relative phase, and every correlation between their
coefficient vectors are allowed. More generally, the parameters and
covariance can be randomized using an arbitrary latent variable
independent of $T$; conditional Gaussianity and the displayed variance
conditions must then hold given that latent variable. Averaging (1)
retains the expected penalties. This permits non-Gaussian unconditional
mixtures, but does not permit an arbitrary conditional cross covariance
between the remote residual and the A/B drivers.

\subsection{A stronger Gaussian adversary}

For a centered Gaussian $R$ with variance $V$ and pseudovariance $P$,
put $p=|P|$. The elementary Gaussian integral gives, for $s>0$,
\[
 \E e^{-s|x+R|^2}
 =D^{-1/2}\exp\left\{-\frac{s x^2(1+sV-s\Re P)}D\right\},
 \qquad D=(1+sV)^2-s^2p^2,
\]
for real $x$. The formula extends to singular covariances by continuity.
The representation
\[
 \sqrt u=\frac1{2\sqrt\pi}\int_0^\infty(1-e^{-su})s^{-3/2}\,ds
\]
shows that, at fixed $V,p$, the norm expectation is largest when $P=-p$:
the larger covariance eigenvalue points perpendicular to the mean.
We allow this choice separately for every $t$.

For this worst orientation the real projection in direction $\theta$
has mean $\ell(t)\cos\theta$ and variance
$[V(t)-p(t)\cos(2\theta)]/2$. Since a Gaussian absolute mean increases
with its variance, we may further enlarge that variance to
\[
 \sigma_\theta(t)^2=\frac12\left[
 g^2+\alpha^2A(t)^2+\beta^2B(t)^2
 +2\alpha\beta|A(t)B(t)|(-\cos2\theta)_+\right].             \tag{2}
\]
The enlargement is valid also for $p<2\alpha\beta|AB|$; simply replacing
$p$ by its maximum in both covariance eigenvalues would \emph{not} be a
valid argument.

Let $F(m,\sigma)=\E|m+\sigma N|$ for a standard real normal $N$.
For each fixed $\theta$, $F(\ell\cos\theta,\sigma_\theta)$ is convex
in the four real fields $(\ell,g,\alpha|A|,\beta B)$. Indeed the
covariance matrix of the last two fields is
$\left(\begin{smallmatrix}1&c\\c&1\end{smallmatrix}\right)$,
where $c=(-\cos2\theta)_+\in[0,1]$, so they can be represented by fixed
correlated real Gaussians. Convexity follows from convexity of the
absolute value. This is the convexity used in the certificate.


\paragraph{The common convex function behind all discretizations.}
Let $\mathcal N(\ell,g,x,y)$ be the supremum of
$\E|\ell+gG+xX+y\overline Y|$ over all jointly proper standard
Gaussians $X,Y$ with $|\E X\overline Y|\le1$, independent of $G$.
Its possible pseudovariances fill exactly the disk $|P|\le2|xy|$;
its variance is $g^2+x^2+y^2$.
For each fixed correlation the norm expectation is convex in the four
real coefficients, hence so is their supremum $\mathcal N$.
The integral $2\E_T\mathcal N(\ell(T),g,\alpha|A(T)|,\beta B(T))$
is therefore a single convex upper bound for every allowed covariance
law. All numerical levels bound this \emph{same} function at vertices.
This is why taking the minimum of bounds from different angular or
radial resolutions remains valid in a parameter-box interpolation.
Convexity of a particular discretization alone would not justify
mixing different discretizations in this way.

\subsection{An exact analytic neighborhood}
We first handle the only region where a numerical upper bound cannot
separate the supremum from its Gaussian value.

\begin{lemma}\label{polar:local}
Theorem~\ref{polar:main} holds whenever
\[
 \frac7{200}\le g^2\le\frac9{200},\qquad
 0\le\alpha,\beta\le\frac1{100}.
\]
\end{lemma}
\begin{proof}
Put $v=g^2$, and let angular brackets denote uniform averaging over
$[0,\pi/2]$. Define
\[
 J_v(t,\theta)=\frac{\sqrt\pi}{2\sqrt v}
      e^{-\ell(t)^2\cos^2\theta/v},
\]
\[
 C_A=\E\langle A^2J_v\rangle,\quad
 C_B=\E\langle B^2J_v\rangle,\quad
 C_\times=\E\langle |AB|(-\cos2\theta)_+J_v\rangle.
\]
The exact checker \texttt{local.cpp} proves, uniformly on the stated
variance interval,
\[
 C_A+\frac25\le\frac1\eps,\qquad
 C_B+\frac9{10}\le\frac1\eps,\qquad C_\times\le\frac8{25}.    \tag{3}
\]
These are elementary positive Gaussian integrals, not assumptions.
Details of their finite enclosure are given below.

Conditionally on $T$, write $R=gG+D$, with $D$ a centered Gaussian
independent of $G$, of variance $\alpha^2A^2+\beta^2B^2$ and
pseudovariance $P$. This decomposition is always possible since
$|P|\le2\alpha\beta|AB|\le\alpha^2A^2+\beta^2B^2$.
For $\psi_g(x)=\E|x+gG|$, the quadratic Taylor term of $2\psi_g$
is $w_v\E|D|^2+k_v\Re\E D^2$, where
\[
 w_v=\langle J_v\rangle,\qquad k_v=\langle J_v\cos2\theta\rangle\le0.
\]
The last sign follows by pairing $\theta$ with $\pi/2-\theta$.
Consequently $|k_v|\le\langle J_v(-\cos2\theta)_+\rangle$,
and (3) bounds the penalized quadratic term by
\[
 -\frac25\alpha^2-\frac9{10}\beta^2+\frac{16}{25}\alpha\beta.   \tag{4}
\]

The fourth derivative of the Gaussian-smoothed planar norm satisfies
\[
 \|D^4\psi_g\|_{\mathrm{op}}\le\frac{3\sqrt\pi}{4g^3}<204.
\]
For completeness, the planar Fourier representation of the norm has
radial density $-(2\pi)^{-1}|\xi|^{-3}$. After convolution with $gG$,
four differentiation factors and polar integration bound a directional
fourth derivative by
\[
 \frac1{2\pi}\left(\int_0^{2\pi}\cos^4\theta\,d\theta\right)
 \left(\int_0^\infty r^2e^{-g^2r^2/4}\,dr\right)
 =\frac{3\sqrt\pi}{4g^3}.
\]
Holder's inequality gives the same bound for four different unit
directions. One can equivalently obtain this representation by
subordinating the Gaussian integral for the norm; the differentiated
integral is absolutely convergent. The inequality $<204$ follows by
squaring the rational bound $\sqrt\pi<1773/1000$ at $g^2=7/200$.

The odd Taylor terms vanish. Moreover
\begin{align*}
 \E|D|^4
 &\le3\E(\alpha^2A^2+\beta^2B^2)^2\\
 &=3(9\alpha^4+14\alpha^2\beta^2+6\beta^4)
 \le27(\alpha^2+\beta^2)^2.
\end{align*}
Thus the fourth-order remainder for $2\psi_g$ is at most
$459(\alpha^2+\beta^2)^2$.
Finally,
\[
 \frac25\alpha^2+\frac9{10}\beta^2-\frac{16}{25}\alpha\beta
 \ge\frac{\alpha^2}{10}+\frac{\beta^2}{4}
                +\frac{\alpha^2+\beta^2}{10},              \tag{5}
\]
since the remaining matrix has diagonal $1/5,11/20$, off-diagonal
$-8/25$, and determinant $19/2500>0$. On the local box,
$459(\alpha^2+\beta^2)\le459/5000<1/10$.
Combining (4)--(5), and using the definition of $S$ for the unperturbed
$gG$, proves the lemma.
\end{proof}

\subsection{How the global certificate works}
All acceptance arithmetic is integer or rational. Floating point is
used only to initialize an integer square-root guess, which is corrected
and verified by exact inequalities before being used.

\subsubsection{Angular enclosure without quadrature assumptions}
For an even positive integer $m$ and every $z\in\mathbb C$,
\[
 |z|\le\tan\frac\pi{2m}\sum_{j=0}^{m-1}
             |\Re(e^{-ij\pi/m}z)|.                         \tag{6}
\]
Indeed the minimum over the phase of $z$ of the sum, divided by $|z|$,
is $\cot(\pi/(2m))$, by the finite geometric sum on a phase cell.
Apply (6) after the worst-orientation reduction and then enlarge each
projected variance by (2). Reflection symmetry leaves $m/2+1$ distinct
angles with multiplicities one at the endpoints and two otherwise.
The certificate uses $m=64,128,256,512$, taking the smallest previously
computed certified bound when a finer level is required.

\subsubsection{Folded normal enclosure}
For $s\ge0$,
\[
 F(s,1)=\sqrt{2/\pi}\left[
 1+\sum_{k=1}^\infty\frac{(-1)^{k-1}s^{2k}}{(2k-1)2^k k!}\right].
\]
On the rational grid $s=j/128$, $0\le j\le1024$, the sum through
$k=201$ is an upper bound, and subtracting the magnitude of term 202
is a lower bound. Convexity supplies upper chords between grid points.
For $s\ge8$, $F(s,1)\le s+10^{-13}$, from
$F(s,1)-s\le\sqrt{2/\pi}e^{-s^2/2}$ and $(8/3)^{32}>10^{13}$.
The expectation is increasing separately in its nonnegative mean and
standard deviation; this deals with rounded input intervals.
Machin's formula bounds $\pi$ by rational arctangent sums. Rational
Taylor sums bound all sines and cosines. No special-function library
value enters an acceptance decision.

\subsubsection{Visible integration and parameter boxes}
Set $r=\sqrt T$. On each cell of $[0,8]$, the fields
\[
 r(a_0+b_0r^2),\quad g,\quad\alpha|r^2-1|,\quad\beta r^2/\sqrt2
\]
have cubic Bernstein control values; $r=1$ is a mesh boundary at every
level. Convexity bounds the integrand by the weighted values at these
controls. The weights are exact outward enclosures of
$\int 2re^{-r^2}B_j^3\,dr$, computed by rational exponential Taylor
bounds. There are $64,128,256,512$ radial cells at the four levels.

For $r\ge8$ and $g,\alpha,\beta\le2$, the angular majorant is bounded
by $r^3+4r^2+r+2$. To see this, use the triangle inequality for each
folded normal and $m\tan(\pi/(2m))<7/4<\sqrt\pi$.
The two-sign tail is less than $1600e^{-64}<2^{-46}$.

At each parameter vertex the checker bounds the objective plus
$\alpha^2/10+\beta^2/4$. Its norm part is convex in $(g,\alpha,\beta)$.
On a rectangular parameter box of side lengths $d_g,d_A,d_B$, the
maximum vertex bound plus
\[
 \frac{d_g^2+d_A^2+d_B^2}{4\eps}                           \tag{7}
\]
is an upper bound throughout the box. The smaller actual price
coefficients after retaining the penalties only improve (7).
Starting from $[0,2]^3$, the checker bisects a longest side whenever
neither the numerical bound nor an analytic rule closes the box.

\subsubsection{Unbounded parameters and the benchmark}
Write $V_0=g^2+\alpha^2+\beta^2$ and
$M=\E\ell(T)^2=a_0^2+4a_0b_0+6b_0^2<19/25$.
Cauchy--Schwarz gives
\[
 \Lambda+\alpha^2/10+\beta^2/4
 \le2\sqrt{M+V_0}-(1/\eps-1/4)V_0<13/10<S
 \qquad(V_0\ge2).                                       \tag{8}
\]
The right side is decreasing there, and its value at two is bounded
by rational comparisons after squaring positive quantities.
This also handles every parameter outside the root cube.

The independent Rice-envelope benchmark calculation is regenerated
with the package. Convex lower integration at $v=1/25$ gives
\[
 B_0=\frac{96683367193980}{2^{46}}\le S.                    \tag{9}
\]
The inherited degree-96 Rice enclosure is verified by its rational
differential residual and endpoint inequalities before use. This
calculation is used only as a lower bound on $S$.

\subsubsection{Acceptance record}
The full parameter tree closes with
\[
 \begin{array}{c|r}
 \text{numerical leaves}&631\\
 \text{analytic local leaves}&17\\
 \text{variance-tail leaves}&5\\
 \text{bisections}&652\\
 \text{distinct parameter vertices}&1226
 \end{array}
\]
The largest numerical-leaf upper bound is
$96683367186193/2^{46}<B_0$. The identity
$631+17+5=652+1$ is a useful count check; the executable reconstructs
the entire partition, so the proof does not rely on that count alone.
Analytic leaves are bounded directly by $S$, not by $B_0$.
In particular, the largest numerical-leaf bound must \emph{not} be
reported as a numerical upper bound on the complete Gaussian family.
Together these three types of leaves prove Theorem~\ref{polar:main}.

\subsection{Enclosure of the local integrals}
The interval $v\in[7/200,9/200]$ is divided into eight cells with
endpoints $(28+k)/800$. There are 2048 radial intervals of length
$1/256$ on $[0,8]$ and 128 angular intervals on $[0,\pi/2]$.
In each cell, take the lower radial endpoint and upper angular endpoint
in the exponential of $J_v$, the upper variance endpoint there, and
the lower variance endpoint in its prefactor. Bound $A^2,B^2,|AB|$
by their elementary polynomial maxima. The exact radial masses are
$e^{-r_l^2}-e^{-r_h^2}$. The angular masses are exactly $1/128$.

For the remaining exponential, a rationally generated table of
$e^{-j/4096}$ uses outward rounding on a $2^{-192}$ lattice and is
finally rounded outwards to the $2^{-46}$ acceptance lattice. Lookup
at the lower grid point is an upper bound. Arguments above 64 use
one acceptance-lattice unit. Tails of all three local integrals are
less than $5(64^2+2\cdot64+2)e^{-64}<2^{-46}$.
The largest returned upper bounds across the eight variance cells are,
for orientation only,
\[
 C_A<0.907744,\qquad C_B<0.417183,\qquad C_\times<0.317200.
\]
Acceptance uses (3) in exact arithmetic, rather than these decimals.

\subsection{Simultaneous improper even channels}
The same method closes an enlarged Gaussian model of all the
zero-diagonal sector covariances, including odd/even interaction.

\begin{theorem}[Gaussian A/B/Y/W comparison]\label{polar:all}
Let $g,\alpha,\beta,\gamma\ge0$, and suppose that $R\mid T=t$ is
centered Gaussian with
\[
 V(t)=g^2+\alpha^2A(t)^2+\beta^2B(t)^2+\gamma^2t,
 \qquad |P(t)|\le2\alpha\beta|A(t)B(t)|+\gamma^2t.
\]
Then
\[
 \boxed{2\E|\ell(T)+R|-\frac{g^2+\alpha^2+\beta^2+\gamma^2}{\eps}
 \le S-\frac{\alpha^2}{10}-\frac{\beta^2}{4}-\frac{\gamma^2}{25}.} \tag{12}
\]
The covariance orientation is unrestricted and may depend on $t$.
In particular the even Gaussian channel need not be proper.
\end{theorem}
\begin{proof}
Enlarge projected variances to
\[
 \frac12\left[g^2+\alpha^2A^2+\beta^2B^2+\gamma^2t
 +(2\alpha\beta|AB|+\gamma^2t)(-\cos2\theta)_+\right].       \tag{13}
\]
The common convex covariance supremum now also allows an independent
unit-variance Gaussian $H$ with arbitrary pseudovariance in the unit
disk, multiplied by the fourth field $\gamma\sqrt t$.
The two disks of pseudovariances add to exactly the disk in the
hypothesis. Thus the convexity and vertex argument remain valid.
The linear field $\gamma r$ uses its own exact Bernstein controls
$r_j=r_l+j(r_h-r_l)/3$; it is not obtained by square-rooting the
controls of $r^2$.

In the local argument the additional quadratic cost is at least
$\gamma^2/10$. Indeed \texttt{local\_even.cpp} certifies
\[
 C_H:=\E\langle T(1+(-\cos2\theta)_+)J_v\rangle,
 \qquad C_H+1/10\le1/\eps
 \quad (7/200\le v\le9/200).
\]
Its largest returned upper value is below $1.214750$; acceptance
again uses rational inequalities. Put $s=\alpha^2+\beta^2$ and
$h=\gamma^2$. The fourth-order remainder is at most
\[
 459s^2+306sh+102h^2,
\]
using $\E A^2T=\E B^2T=3$ and $\E T^2=2$ in the Gaussian
fourth-moment calculation. Therefore the local analytic rule is
\[
 7/200\le g^2\le9/200,\qquad
 459s+306h\le1/10,\qquad102h\le3/50.                      \tag{14}
\]
The remainder in this region is at most $s/10+3h/50$, exactly the
budget left by (5) and the even-channel quadratic penalty after
retaining the three penalties in (12).

For the exterior, use the same fixed-point enclosure on $[0,2]^4$.
The parameter-box correction is the sum of all four squared widths
divided by $4\eps$. With $V_0=g^2+\alpha^2+\beta^2+\gamma^2$, the
variance-tail rule (8) is unchanged, since every retained price is at
most $V_0/4$. The visible tail of the angular majorant is bounded by
$r^3+4r^2+4r+2$, so its two-sign integral is below
$1700e^{-64}<2^{-46}$.

The complete four-parameter cover passes with 4175 numerical leaves,
66 analytic leaves, 15 variance-tail leaves, 4255 bisections, and
9685 distinct vertices. The largest numerical upper bound is
$96683362167072/2^{46}<B_0$. The identity
$4175+66+15=4255+1$ agrees with the reconstructed complete tree.
This proves (12) without numerical slack.
\end{proof}

\subsection{Consequences for the remaining conjecture}
For a general remote cubic $q\in H_{2,1}$, first-chaos $U\in H_{1,0}$
and $V\in H_{0,1}$, set
\[
 R_T=q+A(T)U+B(T)V.
\]
Hermite orthogonality and the charges give
\[
 \E_w|R_T|^2=\|q\|_2^2+A(T)^2\|U\|_2^2+B(T)^2\|V\|_2^2,
 \qquad \E_w R_T^2=2A(T)B(T)\E(UV).
\]
Let $\Gamma_T$ be the centered Gaussian with precisely these two
conditional moments. Define the \emph{signed, visibly averaged error}
\[
 \mathcal D=\E_T\left[
  \E_w(|\ell(T)+R_T|+|\ell(T)-R_T|)
                   -2\E|\ell(T)+\Gamma_T|\right].
\]
The new theorem gives the rigorous implication
\[
 \Lambda(R)\le S-\|U\|_2^2/10-\|V\|_2^2/4+\mathcal D.     \tag{10}
\]
Thus a sufficient remaining target for the odd-sector version of
K$_\eta$ is
\[
 \boxed{\quad\mathcal D\le\|U\|_2^2/10+\|V\|_2^2/4+\eta.\quad} \tag{11}
\]
Unlike fibrewise domination, (11) allows positive non-Gaussian gains,
keeps the visible average, and provides a variance-dependent allowance.
It is \emph{not proved here}. Nor does (10) cover simultaneous $Y/W$
or deterministic diagonal terms. The matched conditional moments alone
do not control $\mathcal D$.

Finally, the theorem transports to profiles
$|a-a_0|,|b-b_0|\le10^{-5}$, with $\eps$ fixed, with additional error
\[
 4\E|Z|\,|a-a_0|+4\E|Z|^3|b-b_0|
 \le\frac{\sqrt\pi}{20000}<\frac{1773}{20000000}=0.00008865.
\]
This follows by applying the norm Lipschitz bound both to the payoff
and to the Gaussian benchmark. It is uniformity for this proved
Gaussian class, not a complete row envelope for arbitrary cubic laws.

\subsubsection{The full zero-diagonal non-Gaussian error}
For $q\in H_{2,1}$, $U\in H_{1,0}$, $V\in H_{0,1}$,
$Y\in H_{1,1}$ and $W\in H_{2,0}$, all on the remote Gaussian space,
put
\[
 R_T=q+A(T)U+B(T)V+\sqrt T(Y+W).
\]
Hermite orthogonality and the circle charges give exactly
\[
 \E|R_T|^2=\|q\|_2^2+A^2\|U\|_2^2+B^2\|V\|_2^2
                  +T(\|Y\|_2^2+\|W\|_2^2),
 \qquad \E R_T^2=2AB\E(UV)+T\E Y^2.
\]
Thus its matched centered Gaussian $\Gamma_T$ satisfies the hypotheses
of Theorem~\ref{polar:all}. With $\mathcal D$ defined as before for
this full $R_T$, we obtain
\[
 \Lambda(R)\le S-\frac{\|U\|_2^2}{10}-\frac{\|V\|_2^2}{4}
 -\frac{\|Y\|_2^2+\|W\|_2^2}{25}+\mathcal D.              \tag{15}
\]
It therefore suffices to bound the signed visibly averaged
non-Gaussian error by the three retained penalties plus $\eta$.
This statement is weaker than fibrewise Gaussian domination and
makes the allowed positive error explicit. Its general validity is
still open. A nonzero deterministic diagonal coefficient changes the
conditional mean and remains outside (15).

\subsection{A quantitative variant for genuinely non-Gaussian shared cubics}
There is also a rigorous, though currently weak, way of using the
Gaussian theorem without assuming Gaussianity of the actual residual.
It approximates the residual \emph{jointly with its coupling drivers}.

Let $q\in H_{2,1}$, $p=\E|q|^2>0$, and
$\kappa=\E|q|^4-2p^2$. Let $U\in H_{1,0}$ and $V\in H_{0,1}$,
with $\alpha=\|U\|_2$, $\beta=\|V\|_2$. Choose orthonormal complex
Gaussian coordinates $X_1,\ldots,X_m$ spanning $U$ and $\overline V$;
$m\le2$. Write $X_j=(x_{2j-1}+ix_{2j})/\sqrt2$ in standard real
Gaussian coordinates, and define
\[
 J=\sum_{j=1}^{2m}\E|\partial_{x_j}q|^2.
\]
No independence of $q$ from these drivers is assumed.

\begin{theorem}[Joint replacement with explicit coupling cost]\label{polar:joint}
For the actual cubic $R_T=q+A(T)U+B(T)V$,
\[
 \Lambda(R)\le S-\frac{\alpha^2}{10}-\frac{\beta^2}{4}
 +2\sqrt{\bigl[p+(\alpha+\beta)^2\bigr]
             \left[\frac{2\kappa}{3p^2}+\frac{8J}{9p}\right]}.
                                                               \tag{16}
\]
In particular the inequality K$_\eta$ holds whenever the last error
is at most $\alpha^2/10+\beta^2/4+\eta$.
\end{theorem}
\begin{proof}
Apply the multivariate Gaussian Stein estimate to the real vector
consisting of $q/\sqrt p$ and $X_1,\ldots,X_m$. Its covariance is
$C=I/2$, by Hermite orthogonality. The estimate of
Nourdin--Peccati--R\'eveillac, Theorem 3.5 and Lemma 3.3
(\url{https://arxiv.org/abs/0804.1889}), bounds Wasserstein distance
by $\sqrt2$ times the root mean square Hilbert--Schmidt Stein-matrix
error. Only the symmetric part of that matrix pairs with a Hessian,
so its symmetric part can be used in the proof of that estimate.

Here the cubic inverse-generator factor is $1/3$, and the linear
factor is one. For the cubic $2\times2$ block, the standard circular
cubic estimate is
\[
 \E\|\tfrac13(\langle Dq_i,Dq_j\rangle)_{i,j}-pI/2\|_{HS}^2
 \le\kappa/3.                                             \tag{17}
\]
One direct verification uses the real cubic identities
$\operatorname{Var}\|DF\|^2=2592A+1944B$ and
$\kappa_4(F)=1296A+1944B$, where $A,B$ are squared tensor
contractions. Thus $\operatorname{Var}\|DF\|^2\le2\kappa_4(F)$.
Circularity gives $\kappa_4(\Re(e^{i\theta}q))=3\kappa/8$;
averaging the gradient-variance identity over $\theta$ proves (17).

A cross entry in the symmetric Stein matrix is
$(2/3)\partial_{x_j}q_i/\sqrt{2p}$, and its transpose has the same
value. All linear--linear entries equal their covariance exactly.
Consequently
\[
 \E\|\Gamma_{\rm sym}-C\|_{HS}^2
 \le\frac{\kappa}{3p^2}+\frac{4J}{9p}.
\]
The target Gaussian has its first complex component independent of
the retained linear drivers. The visibly averaged norm payoff, as a
function of this normalized vector, has Lipschitz constant at most
\[
 2\E_T\sqrt{p+(|A(T)|\alpha+|B(T)|\beta)^2}
 \le2\sqrt{p+(\alpha+\beta)^2}.
\]
The price is unchanged, since all cross second moments vanish.
Combining the joint estimate with Theorem~\ref{polar:main} proves (16).
Lipschitz tests follow from smooth approximation in the Stein estimate.
\end{proof}

For a concrete finite, non-Gaussian example, let
\[
 q_N=\frac1{5\sqrt N}\sum_{j=1}^N H(X_j),\qquad
 U=X_1/10,\quad V=\overline{X_1}/10,\qquad
 H(X)=X(|X|^2-2)/\sqrt2.
\]
Then $p=1/25$, $\kappa=42p^2/N$, and $J=3p/N$.
Formula (16) becomes
\[
 \Lambda(R)\le S-\frac7{2000}+\sqrt{\frac{736}{75N}}.
\]
The exact Wick Stein block improves this further. In standard real
coordinates the Jacobian of $H$ has radial and tangential eigenvalues
$(3T-2)/2$ and $(T-2)/2$, respectively. Hence
\[
 \E\|\Gamma_H-I/2\|_{HS}^2
 =\frac{\E(3T-2)^4+\E(T-2)^4}{144}-\frac12
 =\frac{1000+8}{144}-\frac12=\frac{13}{2}.
\]
For the normalized independent sum this contribution is $13/(2N)$;
the symmetric cross block with $X_1$ contributes $4/(3N)$.
The same joint argument therefore gives the stronger bound
\[
 \boxed{\Lambda(R)\le S-\frac7{2000}+\sqrt{\frac{376}{75N}}.}
\]
It proves exact K for every $N\ge409252$, with no independent
Gaussian buffer. The integer threshold and the Wick moment calculation
are checked in exact rational arithmetic. This is a genuine finite-cubic
consequence, but its kurtosis is extremely close to two; it does
\emph{not} cover the application's difficult kurtosis range.
The quantitative Stein error in (16) is still too large to claim a
new Grothendieck endpoint.

\paragraph{A general joint-Stein sufficient condition.}
The same mechanism is not confined to linear drivers. Form the real
vector $F$ underlying $(q,U,V,Y,W)$, discard deterministic zero
covariance directions, and let $C>0$ be its covariance. Set
$\Gamma_{ij}=\langle DF_i,-D\mathcal L^{-1}F_j\rangle$, where
$\mathcal L$ is the real Gaussian Ornstein--Uhlenbeck generator.
The target Gaussian has exactly the covariance profiles of (15).
The visibly averaged payoff has Lipschitz constant at most $2\sqrt5$,
since $\E[1+A^2+B^2+2T]=5$. Therefore (15) holds with the explicit
upper bound
\[
 \mathcal D\le2\sqrt5\,\|C^{-1}\|_{\rm op}\|C\|_{\rm op}^{1/2}
       \sqrt{\E\|\Gamma_{\rm sym}-C\|_{HS}^2}.              \tag{18}
\]
All quantities are finite polynomial contraction data; the output
dimension is at most ten, independent of the ambient dimension.
The inverse covariance can make (18) numerically poor. It is a proved
sufficient condition, not a claim that it meets the required error
budget for every cubic.

\subsection{IDE reading note}\label{sec:polar-reading}
\paragraph{Where this leaves K.} \autoref{ch:forest}'s fact (c), the Gaussian-profile layer, was numerical (30 optimiser starts). It is now a theorem for the full covariance class, with
explicit penalties: $\alpha^2/10$ ($A$), $\beta^2/4$ ($B$) and $\gamma^2/25$ ($Y/W$). The first-order constants of \autoref{ch:forest}, about $0.47\alpha^2$, $0.93\beta^2$ and $0.70\gamma^2$
at $v^*$, show these certified penalties are deliberately conservative. Every residual now splits as ``matched Gaussian + signed non-Gaussian gap $D$''. The open statement is
$D\le\alpha^2/10+\beta^2/4+\gamma^2/25+\eta$.
\paragraph{Relation to \autoref{ch:stronger} (heuristic \st{?}; the split below is not derived, referee 20 Oct).} The Ornstein--Uhlenbeck-type interpolation between the actual $G$ and its matched Gaussian writes $D=-\int_0^{\pi/2}\phi'(\theta)\,d\theta$, and
$\phi'$ splits into two pieces. The first is the BNC term, which depends on the law of $R_T$ alone and is zero for Gaussian laws; for remote cubics it is the visible-free inequality PVD. The second
is the visible--residual Malliavin cross term $\E\nabla^2|\cdot|(Y)[DL,DG]$, which is nonzero exactly for sectors whose $Z$-dependence overlaps $D_ZL$. A proof of the target can therefore be
organised as BNC (no visible variable) plus a bound on the cross energy by the certified penalties above.
